\documentclass[11pt,a4paper]{article}
\usepackage[top=1.7cm,bottom=1.5cm,left=1.5cm,right=1.5cm,headsep=8pt,footskip=16pt]{geometry}
\usepackage{enumitem}
\usepackage{amssymb}
\usepackage{array}
\usepackage[table]{xcolor}
\usepackage{fancyhdr}
\usepackage{lastpage}
\usepackage{multicol}
\usepackage{polyglossia}
\setdefaultlanguage[calendar=gregorian, numerals=maghrib]{arabic}
\setotherlanguage{english}
\IfFontExistsTF{Amiri}{%
  \newfontfamily\arabicfont[Script=Arabic, Scale=1.08]{Amiri}
}{%
  \newfontfamily\arabicfont[Script=Arabic, Scale=1.08, Path=../fonts/, Extension=.ttf,
    UprightFont=*-Regular, BoldFont=*-Bold, ItalicFont=*-Italic,
    BoldItalicFont=*-BoldItalic]{Amiri}
}
\IfFontExistsTF{AlBattar}{%
  \newfontfamily\khatt[Script=Arabic, Scale=1.15]{AlBattar}
}{%
  \newfontfamily\khatt[Script=Arabic, Scale=1.15, Path=../fonts/,
    BoldFont=ae_AlBattar.ttf]{ae_AlBattar.ttf}
}
\renewcommand{\headrulewidth}{0.8pt}
\renewcommand{\footrulewidth}{0.8pt}
\fancyhf{}
\rhead{\khatt المدرسة العليا للأساتذة ببوزريعة}
\lhead{\khatt ملحقة بومرداس}
\lfoot{\khatt الأستاذ: د. بن عبد الرزاق يوسف}
\cfoot{\thepage/\pageref{LastPage}}
\pagestyle{fancy}
\parindent 0pt
\setlength{\columnsep}{0.6cm}
\setlength{\columnseprule}{0.4pt}
\setlist[enumerate,1]{label=\textbf{\arabic*.}, leftmargin=*, itemsep=5pt, topsep=2pt}
\begin{document}

\linespread{1.08}\selectfont
\begin{center}
{\LARGE\khatt\textbf{الاختبار التاسع عشر في الذكاء الاصطناعي}}\\[4pt]
{\normalsize وضعيات إشكالية ــ أسئلة متعددة الاختيارات ــ الفصول من 1 إلى 6}
\end{center}
\vspace{-2pt}\rule{\textwidth}{0.6pt}\\[2pt]
\textbf{التعليمة:} يتضمن الاختبار خمس وضعيات إشكالية. اقرأ كل وضعية بعناية، ثم اختر التحليل أو التصرف الأنسب
من بين الاختيارات الخمسة؛ لكل وضعية إجابة صحيحة واحدة فقط.
اكتب \textbf{حرف} الإجابة الصحيحة في جدول الإجابات في آخر الاختبار.
لكل وضعية 4 نقاط (المجموع 20 نقطة).
\vspace{4pt}

\noindent
{\large\khatt\textbf{الوضعية 1}}\hfill{\small(الفصل 1)}\par\nopagebreak\vspace{3pt}
\setlength{\fboxsep}{7pt}\colorbox{gray!10}{\parbox{\dimexpr\linewidth-14pt}{صار أستاذ يعتمد على أداة توليدية في إعداد كل دروسه، وينقل مخرجاتها كما هي إلى القسم. لاحظ المفتش أن الدروس منظمة، لكنها لا تراعي أخطاء التلاميذ المتكررة ولا مكتسباتهم السابقة. ويرى الأستاذ أن «الأداة أدرى بالتدريس». ويستعد المفتش لكتابة تقرير عن زيارته، ويريد أن يقترح على الأستاذ طريقة عمل أنسب. ما التقدير الأدق؟}}\par\nopagebreak
\begin{itemize}[leftmargin=1.0cm, labelsep=6pt, itemsep=3pt, topsep=4pt, parsep=0pt]
  \item[\textbf{أ)}] الأستاذ محق، لأن الأداة اطلعت على كل طرائق التدريس، فهي أعرف منه بحاجات تلاميذه، وملاحظة المفتش شكلية لا تأثير لها في جودة التعلّم داخل القسم.
  \item[\textbf{ب)}] يكفي أن يذكر اسم القسم والمستوى في كل طلب لتنتج الأداة دروسًا مكيّفة تمامًا مع أخطاء تلاميذه ومكتسباتهم، فيستمر في نقل المخرجات كما هي دون أي تكييف، فالأداة تعرف خصائص كل قسم من اسمه.
  \item[\textbf{ج)}] المشكلة أن الأداة غير متخصصة في التعليم، ولو استعمل أداة متخصصة لجاءت الدروس مراعية لأخطاء تلاميذه ومكتسباتهم تلقائيًا دون أن يزوّدها بأي معلومة عنهم، وهي متوفرة اليوم في السوق.
  \item[\textbf{د)}] يتوقف الأستاذ عن استعمال الأداة نهائيًا، لأن أي درس تشارك فيه الأداة يكون عامًا بالضرورة، ولا يمكن تكييفه مع القسم مهما أضاف إليه الأستاذ من معرفته.
  \item[\textbf{هـ)}] الأداة لا تعرف قسمه إلا بما يزوّدها به، فتنتج دروسًا عامة؛ ودوره تكييفها بمعرفته بالمكتسبات والأخطاء والسياق والتحقق منها، لأن الأداة مساعد في التحضير لا بديل عن قراره التربوي.
\end{itemize}
\par\vspace{6pt}
\noindent
{\large\khatt\textbf{الوضعية 2}}\hfill{\small(الفصل 3)}\par\nopagebreak\vspace{3pt}
\setlength{\fboxsep}{7pt}\colorbox{gray!10}{\parbox{\dimexpr\linewidth-14pt}{تريد أستاذة شبكة تقويم لوضعية إدماجية في كتابة رسالة، فطلبت «شبكة تقويم للرسالة»، فجاءت معايير عامة تصلح لأي نص، بلا مؤشرات قابلة للملاحظة ولا توزيع للنقاط. والوضعية تستهدف احترام بنية الرسالة وتوظيف الجملة الاستفهامية. وتريد أن تستعمل الشبكة نفسها مع زملائها في المؤسسة لتوحيد التقويم بين الأقسام. أي طلب أنسب؟}}\par\nopagebreak
\begin{itemize}[leftmargin=1.0cm, labelsep=6pt, itemsep=3pt, topsep=4pt, parsep=0pt]
  \item[\textbf{أ)}] تطلب شبكة عامة لكل أنواع النصوص، لأن الشبكة العامة صالحة لكل الوضعيات، والمؤشرات الخاصة ببنية الرسالة والاستفهام تجعل التقويم ضيقًا لا يعكس كفاءة التلميذ.
  \item[\textbf{ب)}] تطلب من الأداة أن تقوّم إنتاجات التلاميذ مباشرة دون شبكة، لأن الأداة تمتلك شبكة داخلية دقيقة لكل الوضعيات، فلا حاجة إلى صياغة معايير ومؤشرات وتوزيع النقاط، وتطابق المنهاج الرسمي.
  \item[\textbf{ج)}] تذكر نص الوضعية وأهدافها والمستوى، وتطلب شبكة في جدول بمعايير الوجاهة والانسجام وسلامة اللغة والإتقان، لكل معيار مؤشرات قابلة للملاحظة مرتبطة ببنية الرسالة والاستفهام، مع توزيع النقاط.
  \item[\textbf{د)}] تذكر نص الوضعية وأهدافها وتطلب الشبكة، ثم تعتمدها دون تجربتها على نماذج من إنتاج التلاميذ، لأن ذكر الأهداف يضمن وحده أن تكون المؤشرات قابلة للملاحظة وعادلة، فتوفّر الوقت على نفسها وزملائها.
  \item[\textbf{هـ)}] تطلب «شبكة تقويم ممتازة ومفصّلة جدًا»، لأن تشديد الصفات يجعل الأداة تستنتج أهداف الوضعية والمستوى وتوزيع النقاط، دون حاجة إلى ذكر نص الوضعية أو أهدافها.
\end{itemize}
\par\vspace{6pt}
\noindent
{\large\khatt\textbf{الوضعية 3}}\hfill{\small(الفصل 4)}\par\nopagebreak\vspace{3pt}
\setlength{\fboxsep}{7pt}\colorbox{gray!10}{\parbox{\dimexpr\linewidth-14pt}{كُلِّف أستاذ ضمن لجنة بإعداد موضوع امتحان رسمي، فأراد رفع مسودة الموضوع إلى أداة ذكاء اصطناعي عامة لتدقيقها لغويًا واقتراح تحسينات قبل تسليمها. والموضوع سري إلى يوم الامتحان. وهو يرى أن التدقيق الآلي سيحسّن جودة الموضوع، والوقت المتبقي قبل التسليم قصير جدًا، وأعضاء اللجنة منشغلون بمهام أخرى. ما التصرف الأسلم؟}}\par\nopagebreak
\begin{itemize}[leftmargin=1.0cm, labelsep=6pt, itemsep=3pt, topsep=4pt, parsep=0pt]
  \item[\textbf{أ)}] يرفعها ويطلب من الأداة كتابيًا ألا تُظهرها لأحد، لأن هذه التعليمة تُلزم الأداة وتحمي سرية الموضوع، فلا حاجة إلى التدقيق مع أعضاء اللجنة، كما تفعل الأنظمة المحمية في المؤسسات الكبرى.
  \item[\textbf{ب)}] يرفعها دون تردد، لأن الأدوات العامة لا تحتفظ بأي نص، ولا يمكن لأي شخص الاطلاع عليه، فالسرية مضمونة تقنيًا ولا تحتاج إلى أي احتياط.
  \item[\textbf{ج)}] يرفعها بشرط أن يستعمل حسابه الشخصي لا حساب المؤسسة، لأن الحساب الشخصي محمي بكلمة مرور، فيبقى الموضوع سريًا ولا يتحمل الأستاذ أي مسؤولية مهنية.
  \item[\textbf{د)}] لا يرفع المسودة إلى أداة عامة، لأن سرية الموضوع التزام مهني، ولا يضمن مصير ما يُرفع؛ فيدقّقها مع اللجنة، أو يستعين بالأداة في مسائل لغوية عامة دون إدراج نص الموضوع.
  \item[\textbf{هـ)}] يرفعها بعد حذف اسم المادة والمستوى، لأن حذف هذه البيانات يجعل الموضوع غير قابل للتعرف، فيبقى سريًا حتى لو اطلع أحد على النص الكامل للأسئلة، ولو تسرّب جزء منه إلى العلن.
\end{itemize}
\par\vspace{6pt}
\noindent
{\large\khatt\textbf{الوضعية 4}}\hfill{\small(الفصل 5)}\par\nopagebreak\vspace{3pt}
\setlength{\fboxsep}{7pt}\colorbox{gray!10}{\parbox{\dimexpr\linewidth-14pt}{أجرت أستاذة استبيانًا عن عادات المطالعة شارك فيه مئة وعشرون تلميذًا، ورفعت النتائج إلى أداة ذكاء اصطناعي لتحليلها. قدّمت الأداة نسبًا مئوية وتحليلًا مقنعًا، لكن الأستاذة لاحظت أن مجموع نسب أحد الأسئلة يتجاوز 100\%. وستعرض النتائج على مجلس الأساتذة ثم على الأولياء في لقاء رسمي. ما التصرف الأنسب؟}}\par\nopagebreak
\begin{itemize}[leftmargin=1.0cm, labelsep=6pt, itemsep=3pt, topsep=4pt, parsep=0pt]
  \item[\textbf{أ)}] تتحقق من الحسابات بنفسها أو بجدول حسابي، وتراجع هل يسمح السؤال باختيارات متعددة قد تفسّر المجموع، وتصحّح التحليل وفق ذلك، لأن سلاسة التحليل لا تضمن صحة الأرقام التي بُني عليها.
  \item[\textbf{ب)}] تتخلى عن الاستبيان كله، لأن خطأ الأداة يعني أن البيانات نفسها فاسدة، ولا يمكن تحليلها بأي طريقة أخرى، حتى بالجداول الحسابية أو بالحساب اليدوي.
  \item[\textbf{ج)}] تعتمد التحليل كما هو، لأن تجاوز المجموع 100\% خطأ شكلي بسيط لا يؤثر في الاستنتاجات، والأداة أدق من الإنسان في الحسابات الإحصائية مهما كانت البيانات، فتوفر وقتًا ثمينًا قبل الاجتماع.
  \item[\textbf{د)}] تحذف السؤال الذي تجاوز مجموعه 100\% من التقرير، وتعتمد باقي النسب والتحليل دون تحقق، لأن خطأ سؤال واحد لا يدل على احتمال أخطاء في الأسئلة الأخرى، وتكفي بقية الأسئلة لبناء الاستنتاجات.
  \item[\textbf{هـ)}] تطلب من الأداة أن «تجعل المجموع 100\%»، وتعتمد الأرقام الجديدة دون مراجعة، لأن تعديل النسب بما يتناسب مع المجموع يصحح التحليل كله تلقائيًا.
\end{itemize}
\par\vspace{6pt}
\noindent
{\large\khatt\textbf{الوضعية 5}}\hfill{\small(الفصل 6)}\par\nopagebreak\vspace{3pt}
\setlength{\fboxsep}{7pt}\colorbox{gray!10}{\parbox{\dimexpr\linewidth-14pt}{يريد أستاذ عرض فيلم وثائقي فرنسي عن الأمير عبد القادر، مع ترجمة عربية آلية مولَّدة للحوار. لاحظ عند المعاينة أن الترجمة تخلط بين أسماء الأعلام، وتترجم بعض التعابير حرفيًا فتغيّر معناها، وتتأخر عن الصوت أحيانًا. ويريد أن يستعمل الفيلم في حصتين، وأن يبني عليه أسئلة للنقاش. ما الخطة الأنسب؟}}\par\nopagebreak
\begin{itemize}[leftmargin=1.0cm, labelsep=6pt, itemsep=3pt, topsep=4pt, parsep=0pt]
  \item[\textbf{أ)}] يعرض الفيلم كما هو، لأن التلاميذ يفهمون السياق العام، والخلط بين الأعلام والترجمة الحرفية لا يؤثران في الرسالة التاريخية للفيلم، والأفلام الوثائقية لا تُعرض لتفاصيلها بل لروحها العامة.
  \item[\textbf{ب)}] يعاين الفيلم كاملًا، ويصحّح الترجمة في المواضع الخاطئة، خاصة الأعلام والتعابير، أو يختار المقاطع التي راجعها، ويقدّم شرحًا شفهيًا لما لم يُصحَّح، لأن الترجمة الآلية قد تحرف المعنى التاريخي.
  \item[\textbf{ج)}] يعرض الفيلم دون ترجمة، لأن الاستماع إلى الفرنسية وحده كافٍ لتلاميذ المتوسط لفهم الأحداث التاريخية، ولا حاجة إلى أي شرح عربي للمقاطع المعروضة، فيتدربون في الوقت نفسه على اللغة الفرنسية.
  \item[\textbf{د)}] يطلب ترجمة ثانية بأداة أخرى، ويعرض الفيلم بالترجمة الجديدة دون معاينة، لأن الأداة الثانية تصحح أخطاء الأولى بالضرورة في الأعلام والتعابير والتزامن، فيكون العرض جاهزًا في وقت قصير ودون أي جهد.
  \item[\textbf{هـ)}] يتخلى عن الفيلم نهائيًا، لأن أي ترجمة آلية غير قابلة للتصحيح، ولا يمكن الاستفادة من الأفلام الأجنبية في تدريس التاريخ الوطني بأي طريقة.
\end{itemize}
\par\vspace{6pt}
\begin{center}
{\khatt\textbf{جدول الإجابات}}\\[3pt]
\renewcommand{\arraystretch}{1.6}
\begin{tabular}{|>{\centering\arraybackslash}m{1.9cm}|*{5}{>{\centering\arraybackslash}m{1.2cm}|}}
\hline
\textbf{رقم السؤال} & \textbf{1} & \textbf{2} & \textbf{3} & \textbf{4} & \textbf{5} \\ \hline
\textbf{الإجابة} &  &  &  &  &  \\[8pt] \hline
\end{tabular}

\vspace{8pt}
\renewcommand{\arraystretch}{2.2}
\begin{tabular}{|>{\centering\arraybackslash}m{7.5cm}|>{\centering\arraybackslash}m{3.2cm}|>{\centering\arraybackslash}m{3.2cm}|}
\hline
\textbf{الاسم واللقب} & \textbf{الفوج} & \textbf{العلامة} \\ \hline
 & & \hfill /20 \\ \hline
\end{tabular}
\end{center}
\end{document}
